\section{The observer-conditioned deal posterior}
\label{sec:belief}

This section builds the object the rest of the report measures, and then reports
the first of the study's two principal negative results: the exact posterior is a
worse predictor than the approximation the engine actually deploys.

The order matters, because the result is easy to misread in either direction.
Two facts about the exact posterior are both true. It resolves distinctions the
approximation cannot --- it is the exact conditional count under the rules, with
no independence assumption anywhere in it --- and it predicts where cards are
worse than the approximation does. Resolution and calibration are different
properties, and this section keeps them apart throughout. The reason the exact
object is the worse predictor is not that the approximation is secretly better at
counting. It is that the exact object is exact \emph{with respect to a uniform
prior over deals}, and that prior is wrong.

\subsection{The hidden state is the initial deal}
\label{sec:belief-state}

In the dialect of \S\ref{sec:rules} a card changes hands in exactly one way, a
successful ask, which every player observes; declarations remove cards from play,
and that too is public. The v0.4 report proves the consequence
\cite{fishbot04}: a card whose location has never been publicly revealed is still
held by the player to whom it was dealt, the joint hidden state at any time is
exactly the initial deal, and the current holder of every card is a
deterministic function of the initial deal and the public history. The argument
is one line --- every change of holder is an announced event, so a card with no
such event has not moved --- and it is cited here rather than reproved.

What that buys is the elimination of the filtering recursion that dominates
inference in most imperfect-information card games. There is no belief that
decays and no dynamics to approximate, only a fixed hidden object and an
accumulating set of logical constraints on it. Every player computes the
posterior from the same public history but conditions additionally on their own
private hand, so the six posteriors differ and no single distribution is shared.

The constraints are accumulating rather than merely present, and the rate at
which they accumulate governs whether an exact endgame treatment is worth
building. A game reaches a decided state after a mean of \vsixCollapseEvents{}
public events, which is the figure \S\ref{sec:mechanisms} uses to size the
endgame regime.

\paragraph{Notation.} $\mathcal{H}$ is the set of live half-suits and $H \in
\mathcal{H}$ one of them. Fix an observer $i$. $U$ is the set of cards in live
half-suits whose holder is not publicly known to $i$, and for $c \in U$ the map
$\sigma : U \to \{0,\dots,5\}$ gives that card's holder, constant in time.
$M_c \subseteq \{0,\dots,5\}$ is the surviving support of $\sigma(c)$, and $q_p$
is the capacity of seat $p$, the number of cards of $U$ that seat must hold. $Z$
is the number of assignments $\sigma$ consistent with the constraints below, the
partition function of the posterior. The three quantities the policy consumes are
the per-card ownership marginal $\mu_{c,p} = \Pr[\sigma(c) = p]$; the probability
$\tau(H,T)$ that team $T$ owns every card of half-suit $H$; and the probability
$\alpha(A)$ that a named allocation $A$ --- an assignment of all six cards of one
half-suit to named seats --- is correct. Unhatted symbols are values from the
exact engine, hatted ones ($\hat\mu_{c,p}$, $\hat\tau$, $\hat\alpha$) the same
values from the approximate path. Probabilities are conditional on the observer's
information throughout and the conditioning bar is suppressed.

\paragraph{The constraint system.} The observer's information is exactly five
constraints on $\sigma$.

\begin{description}[style=unboxed,leftmargin=0pt,itemsep=4pt,parsep=0pt]
  \item[(C1) Own hand.] $\sigma(c) \ne i$ for every $c \in U$, and every card $i$
        holds is resolved.
  \item[(C2) Public transfers and correct declarations.] Any card whose holder
        has been revealed leaves $U$ with its holder known exactly. A successful
        ask reveals the holder \emph{before} the transfer, and that is the value
        relevant to every earlier constraint.
  \item[(C3) Exclusions.] An ask for $c$ proves the asker did not hold $c$; a
        miss proves the target did not hold $c$. Both are permanent, because
        unresolved cards never move. The surviving support $M_c$ records them.
  \item[(C4) Capacities.] Hand counts are public, so for every player $p$
        \begin{equation}
          \bigl|\{c \in U : \sigma(c) = p\}\bigr| \;=\; q_p \;:=\; h_p - k_p ,
          \label{eq:capacity}
        \end{equation}
        where $h_p$ is $p$'s public card count and $k_p$ the number of live cards
        already known to be $p$'s. Note $\sum_p q_p = |U|$ automatically, and
        that a declaration --- correct or not --- is absorbed by
        \eqref{eq:capacity} without special handling, because the count drop it
        causes is public.
  \item[(C5) Ask legality.] An ask for $c$ in half-suit $H$ by player $a$ proves
        that $a$ held some \emph{other} card of $H$ at that moment. Because
        unresolved cards never move, this is the time-independent disjunction
        \begin{equation}
          \bigvee_{d \in D} \bigl[\sigma(d) = a\bigr],
          \qquad D = \{d \in H \setminus \{c\} : d \text{ unresolved},\; a \in M_d\},
          \label{eq:disj}
        \end{equation}
        vacuous if some card of $H$ is already known to be $a$'s. Every such
        certificate is confined to a single half-suit, which is the structural
        fact \S\ref{sec:belief-count} exploits.
\end{description}

The prior is uniform over deals consistent with the rules, so the posterior is
uniform over the assignments $\sigma$ satisfying (C1)--(C5) and $\Pr[\sigma] =
1/Z$ for each of them. That sentence contains the whole of
\S\ref{sec:belief-exactworse} in advance: the object is exact with respect to
\emph{that} prior, and the prior is a modelling choice, not a rule of the game.

\subsection{The exact count law}
\label{sec:belief-count}

Set (C5) aside. Counting assignments subject to (C3) and (C4) is a
degree-constrained bipartite counting problem, solved exactly by a
forward--backward recursion over the vector $n = (n_0,\dots,n_5)$ of
\emph{remaining} capacities. Ordering $U$ as $c_1,\dots,c_{|U|}$, $F_j(n)$ counts
placements of the first $j$ cards leaving $n$ unused and $B_j(n)$ placements of
the rest out of $n$:
\begin{align}
  F_j(n) &= \sum_{p \,\in\, M_{c_j} :\, n_p < q_p} F_{j-1}(n + e_p),
  &F_0(q) &= 1, \label{eq:fwd}\\
  B_j(n) &= \sum_{p \,\in\, M_{c_{j+1}} :\, n_p \ge 1} B_{j+1}(n - e_p),
  &B_{|U|}(\mathbf{0}) &= 1, \label{eq:bwd}\\
  Z &= F_{|U|}(\mathbf{0}) \;=\; B_0(q), \label{eq:Z}
\end{align}
with $e_p$ the $p$-th unit vector. Every card consumes exactly one unit of
capacity, so the layer index is recoverable from $n$ and the $|U|$-layer table
collapses to two arrays of at most $10^5$ entries (\S\ref{app:inf-count}). These
recursions count rather than sample, so they and the per-card marginal they
contract to are exact for (C1)--(C4).

\paragraph{Ask legality by half-suit block enumeration.} Constraint (C5) is
disjunctive and breaks the product form of \eqref{eq:fwd}--\eqref{eq:bwd}, and
inclusion--exclusion over the tens of certificates a game accumulates would cost
$2^k$ passes. The escape is structural: every certificate lives inside one
half-suit, so $U$ partitions into half-suit blocks admitting one of two exact
treatments. A block with no live certificate is expanded card by card as in
\eqref{eq:fwd}--\eqref{eq:bwd}. A block with a live certificate is enumerated
exhaustively --- at most $\prod_{c \in H}|M_c| \le 5^6$ assignments --- each
tested against \eqref{eq:disj} directly, and the survivors aggregated into a
count-vector table
\begin{equation}
  g_H(t) \;=\; \bigl|\{\text{assignments } \beta \text{ of } H :
    \mathrm{count}(\beta) = t,\; \beta \models \text{certificates of } H\}\bigr| ,
  \label{eq:blocktable}
\end{equation}
where $\mathrm{count}(\beta)_p$ is the number of $H$'s unresolved cards $\beta$
gives to seat $p$. Either way a block consumes a known number of cards, so the
layer identity survives and the dynamic program runs over blocks rather than
cards. Writing $S_t$ for the path mass of count vector $t$ at its block's entry
layer,
\begin{equation}
  S_t \;=\; \sum_{|n| \,=\, \text{entry layer}} F(n)\, B(n - t),
  \label{eq:pathmass}
\end{equation}
the three quantities of \S\ref{sec:belief-state} are read straight off the
tables:
\begin{align}
  \mu_{c,p} &= \frac{1}{Z}\sum_t g_H^{(c \to p)}(t)\, S_t
  &&\text{(per-card ownership)} \label{eq:q1}\\
  \alpha(A) &= \frac{S_{t(A)}}{Z}
  &&\text{(the declaration query)} \label{eq:q2}\\
  \tau(H,T) &= \frac{1}{Z}\sum_{t \,:\, \mathrm{supp}(t) \subseteq T} g_H(t)\, S_t
  &&\text{(does one team hold the half-suit?)} \label{eq:q3}
\end{align}
All three are exact under the uniform prior; no independence assumption and no
sampling enters. \S\ref{app:inf-blocks} gives the mechanics of both block kinds.

\paragraph{Exactness here is cheap.} An exact rebuild costs \vsixCountCost{}~ms
mean over \vsixCountStates{} states, against \vsixSinkhornCost{}~microseconds for
the approximate path. That ratio is large but the absolute cost is small, and it
is small enough that the result of \S\ref{sec:belief-exactworse} is not a budget
story. The exact posterior is not too expensive to deploy. It is worse.

\paragraph{What is verified and what is merely unfalsified.} Two exactness claims
in this machinery are of different standing, and the distinction is stated here
rather than buried in the appendix. The count law and the block enumeration are
checked against a brute-force enumerator that factorises nothing
(\S\ref{app:inf-audit}), and the agreement is at integer rather than rounding
scale. The team-allocation search, by contrast, does not re-check the survival
constraint of \eqref{eq:disj} on the allocation it returns; it has failed zero
times, which is not the same as a proof. \S\ref{app:inf-latent} lists that case
and two others --- a count routine that falls back to a greedy approximation for
heterogeneous non-singleton masks, and an enumeration truncation that leaves a
partial table which is then normalised as if complete --- each with its
reachability argument.

\paragraph{The approximate path.} Six observers invoke the posterior after every
public event, which puts it in the inner loop of every large experiment, so the
engine also implements an approximation and that is what ships. The Fast path
keeps the exact bookkeeping of (C1)--(C3) and the exact capacities (C4), fits
them by Sinkhorn scaling of the support matrix, and interleaves a conditioning
step for (C5) that reweights a certificate's cards by treating their entries as
independent Bernoulli indicators before capacity fitting is re-run
(\S\ref{app:inf-sinkhorn}). The independence assumption is false, because the
capacity constraints that make \eqref{eq:q2} necessary also couple the cards
inside a certificate. Every performance number in this report was produced by
this path.

\subsection{The policy prior}
\label{sec:belief-prior}

The posterior of (C1)--(C5) is exact with respect to the rules and is not the
full Bayesian posterior over deals, which would additionally weight each deal $x$
by the likelihood that the observed actions would have been chosen given $x$.
That likelihood is playstyle-dependent: a player who has taken many turns and
never asked in half-suit $H$ is, under any plausible policy, less likely to hold
cards of $H$, and the rules alone say nothing about this.

The engine implements a two-parameter family of such corrections as prior weights
inside the same machinery,
\begin{equation}
  w(c,p) \;=\; \exp\!\bigl(\theta\, a_{p,H(c)} - \phi\,(a_p - a_{p,H(c)})\bigr),
  \label{eq:prior}
\end{equation}
where $a_{p,H}$ counts $p$'s public asks in half-suit $H$, $a_p$ their total
public asks, and $H(c)$ is the half-suit of card $c$. The weights replace the
uniform initialisation of the Sinkhorn fit, so $\theta = \phi = 0$ recovers the
policy-agnostic posterior of \S\ref{sec:belief-state}. Both parameters are
fitted jointly with the rest of the policy (\S\ref{sec:fitting}); the shipped
values are $\theta = \vsixPriorTheta{}$ and $\phi = \vsixPriorPhi{}$. Deleting
the prior from the shipped configuration, leaving every other coordinate at its
fitted value, costs \vsixPriorDeleteCost{} points, interval
[\vsixPriorDeleteCI{}].

Two architectural facts follow, and they are what makes
\S\ref{sec:belief-exactworse} possible. \emph{First}, \eqref{eq:prior} is the
only channel by which behaviour enters the belief; everything else in
(C1)--(C5) is a rule. \emph{Second}, the approximate path has that channel and
the exact path structurally cannot: the block construction of
\S\ref{sec:belief-count} takes the deduction state alone and has no parameter
through which an opponent model could reach it. The two paths are therefore not
a fast approximation and its exact reference. They are two different posteriors
over the same hidden state, one policy-agnostic and exact under a uniform deal
prior, the other policy-aware and approximate.

\subsection{The exact posterior is a worse predictor than the approximation}
\label{sec:belief-exactworse}

Both posteriors were scored as predictors, on the same states, by two measures
that do not involve playing a game: predictive log loss over the unresolved cards
of each state, and the realised hit rate of each posterior's own argmax ask ---
the ask the posterior itself rates most likely to succeed, scored against the
truth. The probe is \texttt{fish v6probe --mode=belief}
(\S\ref{sec:engine-impl}), and the source of record is
\filepath{research/v06/notes/R12-mechanism-trials.md}, \S3.

The comparison carries no fitted weights on either side. It scores posteriors,
not policies, so nothing in it depends on which posterior the ask weights were
tuned against. That is why it can settle a question that a substitution
experiment in play cannot.

\begin{table}[ht]
\centering
\caption{Each posterior scored as a predictor on identical states: mean
predictive negative log likelihood over the unresolved cards, the probability
each posterior assigns to its own argmax ask, and the realised hit rate of that
ask. Rows are the exact count law under the uniform deal prior, the approximate
path with the policy prior deleted, and the approximate path at the shipped
prior, followed by the $\theta$ sweep. Lower NLL and higher hit rate are better.
The unit is one card for the NLL column and one decision for the hit column.}
\label{tab:belief}
\small
%
\begin{tabular}{lrrrr}
\toprule
Posterior & Cards & Mean NLL & Argmax $p$ & Argmax hit \\
\midrule
exact (uniform prior) & 140,661 & 1.42246 & 0.4660 & 47.94 \\
sinkhorn th=0.000 ph=0.122 & 140,661 & 1.39083 & 0.4896 & 49.99 \\
sinkhorn th=0.200 ph=0.122 & 140,661 & 1.38663 & 0.4939 & 50.75 \\
sinkhorn th=0.300 ph=0.122 & 140,661 & 1.38465 & 0.4972 & 51.06 \\
sinkhorn th=0.445 ph=0.122 & 140,661 & 1.38218 & 0.5035 & 51.49 \\
sinkhorn th=0.600 ph=0.122 & 140,661 & 1.38055 & 0.5113 & 51.59 \\
sinkhorn th=0.800 ph=0.122 & 140,661 & 1.38055 & 0.5211 & 51.64 \\
sinkhorn th=1.100 ph=0.122 & 140,661 & 1.38311 & 0.5336 & 51.64 \\
sinkhorn th=1.500 ph=0.122 & 140,661 & 1.38758 & 0.5459 & 51.51 \\
sinkhorn th=2.000 ph=0.122 & 140,661 & 1.39221 & 0.5555 & 51.40 \\
\bottomrule
\end{tabular}

\end{table}

The ordering is unambiguous and it is the wrong way round. The exact count law
under a uniform deal prior scores \vsixBeliefExactNLL{} nats per card and
\vsixBeliefExactHit{}\% argmax hit rate. The approximate path with only the
certificate half of the prior deleted ($\theta = 0$) scores \vsixBeliefNoPriorNLL{} and
\vsixBeliefNoPriorHit{}\%; with the prior deleted outright ($\theta = \phi = 0$,
\filepath{research/v06/results/F2-belief-noprior.txt}) it scores $1.39339$ and $49.14\%$ --- still
ahead of the exact law on both measures. The shipped prior scores \vsixBeliefShippedNLL{} and
\vsixBeliefShippedHit{}\%. The exact object is the worst predictor of the three,
and the gap between the second and third rows is the policy prior alone:
\numThetaPriorHitGain{} points of argmax hit rate and \numThetaPriorNatGain{}
nats.

The same ordering appears on the full live candidate set rather than on the card
marginals. Over the ask decisions of a \vfive{} mirror, the approximate argmax
converts \vsixFullFastFive{}\% of the time and the exact argmax
\vsixFullExactFive{}\%, on identical states, with the two argmaxes disagreeing at
\vsixDisagreeFive{}\% of decisions; hindsight converts
\vsixFullBestFive{}\%. Substituting the exact posterior into the shipped policy
therefore does not merely fail to help. It replaces the better predictor with the
worse one at every decision on which they differ.

Three consequences follow, and the note draws all three.

\begin{enumerate}
  \item \textbf{``Exact'' is exact under a uniform-deal prior, and that prior is
        wrong.} The exact count law is the worst predictor in
        Table~\ref{tab:belief}. Part of the gap is \eqref{eq:prior}: deleting
        the prior costs the approximation about a point and a half of argmax hit
        rate. But it is not the whole gap --- with both prior parameters deleted
        the approximation still predicts better than the exact law, so the
        remainder is the Sinkhorn fit's own maximum-entropy smoothing, which
        hedges better than an exact conditional distribution taken under a prior
        that is wrong. Exactness with respect to the
        rules is not exactness with respect to the game, because the observed
        actions carry information the rules do not encode. Calling the block
        posterior ``the exact posterior'' without that qualification --- as the
        earlier reports in this corpus do --- names it after a property it has
        only relative to an assumption nobody tested.
  \item \textbf{The residual belief error is not an approximation error to be
        engineered away.} This settles the corpus's standing open question
        --- whether a matched-budget refit of the ask weights under
        \texttt{belief=block} would recover the deficit the v0.4 substitution
        ablation measured --- in the negative. The refit is not worth running,
        because the object it would fit is a strictly less informative posterior.
        The question stayed open for two studies because it was posed as a
        budget question, and it is not one: \S\ref{sec:belief-count} shows the
        exact path is affordable, and this subsection shows it is worse anyway.
  \item \textbf{The predictive optimum is not the playing optimum.} Raising
        $\theta$ improves prediction well past its fitted value: the sweep in
        Table~\ref{tab:belief} minimises NLL at $\theta \approx
        \numThetaPredictiveOpt{}$, against the shipped $\numThetaShipped{}$. In
        play the same change loses --- \numThetaSixtyCost{} points of win rate at
        $\theta = \numThetaTrialLow{}$ and \numThetaElevenCost{} at $\theta =
        \numThetaTrialHigh{}$, at \numThetaTrialN{} games per cell. The ask weights were fitted at the
        shipped $\theta$ and do not transfer to a different one. Belief and
        policy are therefore not separable objects to be optimised in turn; any
        change to $\theta$ has to be refitted jointly with the weights that
        consume it, which is what the fit of \S\ref{sec:fitting} does.
\end{enumerate}

The third consequence is the one that generalises beyond this game. A better
posterior is not a better agent, and an inference component evaluated on
predictive loss can be improved into a policy that plays worse. The two
objectives coincide only when the downstream policy is optimal for whatever
posterior it is handed, and no fitted linear policy is.
